\scanpage{66}
Cette interprétation montre que la catégorie \inserted{isotopique} des cartes (triangulées) pondérées (de type $\{0,1,\infty\}$) est équivalente à \inserted{(une ss-cat. pleine de)} celle des \inserted{$\{1,\tau\}$-}revêtements \inserted{finis} étales de $U_{0,3}$, \deleted{tels que}\inserted{disons ceux} satisfaisant à une condition supplémentaire (savoir que sur l'ens. à op. correspondant, $\tau_0,\tau_1,\tau_\infty$ opèrent sans pts fixes). Mais on vérifie, $\unsure{\Longleftrightarrow}$ de suite que \deleted{\unsure{pour le}} pour un $\{1,\tau\}$ revêtement de $U_{0,3}$, les points fixes des $\tau_i$ correspondants correspondent aux composantes connexes de $U^\tau$ qui sont au dessus de l'arc de $U_{0,3}(\mathbb R)$ opposé à $P_i$ i.e \deleted{\unsure{au b}} de centre $Q_i$ -- donc la condition signifie que \deleted{$\unsure{U^\tau=\varnothing}$} l'on a
\[
\text{(25)}\qquad U^{\prime\tau}=\varnothing\quad\text{i.e.}\quad X^{\prime\tau}=\varnothing
\]
où $X'$ est le revêtement ramifié de $\mathbb P^1(\mathbb C)$ qui prolonge $U'$. [NB \deleted{\unsure{si on avait}} a priori $X^{\prime\tau}$ est une courbe analytique réelle \deleted{plongée}\inserted{contenue} dans $X'$; \deleted{\unsure{qui} $X^{\prime\tau}$} est \inserted{donc} vide dès que sa trace sur $U'=X'\smallsetminus\unsure{\text{ens fini}}$ \deleted{(ens fini)} l'est.)] \inserted{NB.} En termes de la carte pondérée initiale, $X'$ est le revêtement de ses orientations \inserted{locales}, $\tau_{X'}$ \deleted{\unsure{en}} l'autom. \unsure{d'ordre} \unsure{passant} à l'orientation opposée..., qui en effet doit être sans points fixes).

Indépendamment de la condition restrictive (25), la donnée du $\{1,\tau\}$ revêtement \inserted{ramifié} de $\unsure{\boldsymbol{\Sigma}|\Sigma}=\mathbb P^1(\mathbb C)$,
\margin{65}
